\section*{Bab \ref{ch:b1:weathering-clays} --- Pelapukan dan lempung}

\begin{solution}{exo:b1:weathering-clays:1}
\ce{K+}: $1/1.38 = 0.72$; \ce{Fe^2+}: $2/0.78 = 2.56$; keduanya di bawah 3, kation
larut. \ce{Fe^3+}: $3/0.645 = 4.65$; \ce{Ti^4+}: $4/0.605 = 6.61$; keduanya
di antara 3 dan 10, hidrolisat. Di tanah yang tergenang air, tanpa oksigen, besi
dilepaskan sebagai \ce{Fe^2+}, yang tetap dalam larutan dan terbawa oleh air; di
tanah yang berdrainase baik, oksigen mengubahnya menjadi \ce{Fe^3+}, yang langsung
mengendap sebagai oksida dan hidroksida dan tinggal.
\end{solution}

\begin{solution}{exo:b1:weathering-clays:2}
A mengembang dengan glikol dan menyusut ketika dipanaskan: sebuah \omterm{def:b1:weathering-clays:clay-mineral}{smektit}. B
berada pada \SI{7}{\angstrom} dan hancur pada \SI{550}{\celsius}: \omterm{def:b1:weathering-clays:clay-mineral}{kaolinit}. C,
pada \SI{10}{\angstrom} apa pun perlakuannya: \omterm{def:b1:weathering-clays:clay-mineral}{illit}. D, pada \SI{14}{\angstrom}
tanpa mengembang atau menyusut: \omterm{def:b1:weathering-clays:clay-mineral}{klorit}.
\end{solution}

\begin{solution}{exo:b1:weathering-clays:3}
Dua ortoklas memberikan 2 aluminium satu \omterm{def:b1:weathering-clays:clay-mineral}{kaolinit}; $6 - 2 = 4$ silikon keluar
sebagai \ce{H4SiO4}; 2 \ce{K+} keluar, diseimbangkan oleh 2 \ce{H+}; oksigen,
$16 + w = 9 + 16$, memberi $w = 9$:
\[
  \ce{2 KAlSi3O8 + 2 H+ + 9 H2O -> Al2Si2O5(OH)4 + 2 K+ + 4 H4SiO4}.
\]
Hidrogen cocok: $2 + 18 = 20 = 4 + 16$. Seluruh kalium keluar, dan $4/6 =
2/3$ dari silikon.
\end{solution}

\begin{solution}{exo:b1:weathering-clays:4}
(a) Massa tidak berubah: $10 \times 999.87/916.72 = \qty{10.91}{cm^3}$.
(b) Antara sekitar \qty{-4}{\celsius} dan \qty{-15}{\celsius}. Tepat di bawah nol
es tidak dapat membangun tekanan yang cukup untuk membaji retakan; pada beku yang
lebih keras air tidak lagi bergerak melalui pori-pori menuju es, dan es berhenti
tumbuh. (c) Dengan
\cref{prop:b1:weathering-clays:surface} dengan $n = \qty{1}{m}/\qty{1}{cm} = 100$,
luasnya dikalikan 100: dari \qty{6}{m^2} menjadi \qty{600}{m^2}.
\end{solution}

\begin{solution}{exo:b1:weathering-clays:5}
$S/m = 6/(\rho d)$: untuk \qty{1}{cm}, $6/(2600\times0.01) =
\qty{0.231}{m^2.kg^{-1}}$; untuk \qty{1}{mm}, \qty{2.31}{m^2.kg^{-1}}; untuk
\qty{1}{\micro m}, \qty{2308}{m^2.kg^{-1}}, atau \qty{2.3}{m^2.g^{-1}}.
Reaksi kimia berlangsung di permukaan butir, sehingga lajunya per kilogram batuan
tumbuh seiring permukaan spesifiknya, sepuluh ribu kali dari sentimeter ke
mikrometer. \omterm{def:b1:weathering-clays:weathering}{Pelapukan fisik} memotong batuan menjadi potongan yang makin kecil dan
membuka retakan tempat air masuk: ia melipatgandakan permukaan tempat \omterm{def:b1:weathering-clays:weathering}{pelapukan
kimia} dapat bekerja.
\end{solution}

\begin{solution}{exo:b1:weathering-clays:6}
Satu ton albit adalah $\num{e6}/262.22 = \qty{3813.6}{mol}$. Menurut
\cref{prop:b1:weathering-clays:feldspars}, dua mol memberi satu \omterm{def:b1:weathering-clays:clay-mineral}{kaolinit}:
$\qty{1906.8}{mol} \times \qty{258.16}{g.mol^{-1}} = \qty{492.3}{kg}$. Dua dari
tiga atom silikon keluar: $2\times3813.6 = \qty{7627.2}{mol}$ dari
\ce{SiO2}, \qty{458.2}{kg}. Seluruh natrium keluar: $3813.6\times22.99 =
\qty{87.7}{kg}$. Kurang dari setengah massa \omterm{def:b1:rock-forming-minerals:groups}{feldspar} tersisa sebagai lempung.
\end{solution}

\begin{solution}{exo:b1:weathering-clays:7}
Dengan \cref{prop:b1:weathering-clays:immobile}, satu kilogram granit menyisakan
$0.30/0.60 = \qty{0.5}{kg}$ saprolit. Fraksi yang hilang: natrium $1 -
(0.35/3.5)(0.5) = 0.95$; silikon $1 - (60/72)(0.5) = 0.583$; aluminium $1 -
(27/14)(0.5) = 0.036$. Aluminium praktis tetap, natrium hampir seluruhnya hilang,
dan hampir \qty{60}{\percent} silikon juga, meskipun sebagian besar yang tersisa
tertahan dalam kuarsa: \omterm{def:b1:rock-forming-minerals:groups}{feldsparnya} sedang berubah menjadi \omterm{def:b1:weathering-clays:clay-mineral}{kaolinit}, \omterm{def:b1:weathering-clays:hydrolysis-degree}{monosiallitisasi}.
\end{solution}

\begin{solution}{exo:b1:weathering-clays:8}
Pirofilit: $x = y = 0$, muatan 0. Lapisan muskovit: $x = 1$, muatan $-1$.
Lapisan beidelit: $x = 0.5$, muatan $-0.5$. Lapisan montmorilonit: $y = 0.33$,
muatan $-0.33$. Kedua lapisan \omterm{def:b1:weathering-clays:clay-mineral}{smektit}, beidelit dan montmorilonit, dapat
mengembang: muatan kecil mereka diseimbangkan oleh sedikit kation antarlapisan,
yang memegang lapisan-lapisannya longgar dan menarik air ketika terhidrasi.
Pirofilit tidak memiliki kation antarlapisan untuk dihidrasi, dan muatan $-1$
muskovit diseimbangkan oleh ion kalium yang mengunci lapisan-lapisannya.
\end{solution}

\begin{solution}{exo:b1:weathering-clays:9}
(a) Mg: 2 dan 2; Si: 1 dan 1; C: 4 dan 4; O: $4 + 8 + 4 = 16$ dan $12 + 4 =
16$; H: 8 dan $4 + 4 = 8$; muatan: 0 dan $4 - 4 = 0$. (b) Massa molar
forsterit adalah $2\times24.305 + 28.085 + 4\times15.999 =
\qty{140.69}{g.mol^{-1}}$, sehingga satu ton adalah \qty{7108}{mol}. Ia mengonsumsi
$4\times7108 = \qty{28431}{mol}$ dari \ce{CO2}, \qty{1251}{kg}; dengan dua
magnesium per satuan rumus, \cref{prop:b1:weathering-clays:co2} menyimpan
$2\times7108$ mol, \qty{626}{kg}, setengah dari jumlah yang dikonsumsi. (c) Satu ton
kalsit, \qty{9991}{mol}, mengonsumsi \qty{9991}{mol} \ce{CO2} (\qty{440}{kg})
saat melarut, dan mengembalikan semuanya ketika karbonat mengendap kembali: tidak
ada yang tersimpan.
\end{solution}

\begin{solution}{exo:b1:weathering-clays:10}
(a) $3 < 3/r < 10$ memberi $\qty{0.3}{\angstrom} < r < \qty{1.0}{\angstrom}$.
(b) $3/0.27 = 11.1 > 10$: boron bergerak sebagai oksianion, borat. (c) $2/0.83 =
2.41$, kation larut; $4/0.53 = 7.55$, hidrolisat. Dalam air yang miskin oksigen
mangan larut sebagai \ce{Mn^2+}; di tempat mata air bertemu udara, ia teroksidasi
menjadi \ce{Mn^4+}, yang langsung mengendap sebagai oksida hitam pada batuan. (d)
\omterm{def:b1:weathering-clays:ionic-potential}{Potensial ioniknya}, $4/0.40 = 10.0$, berada pada batas antara hidrolisat dan
oksianion: silikon tidak ditahan seperti aluminium maupun dibawa pergi seperti
sulfur, dan fraksi yang tercuci bergantung pada seberapa banyak air yang lewat,
yang menjadi dasar tiga derajat \omterm{def:b1:weathering-clays:hydrolysis}{hidrolisis}.
\end{solution}

\begin{solution}{exo:b1:weathering-clays:11}
(a) Dengan \qty{12.011}{g.mol^{-1}} karbon, \num{11.7e12} hingga
\qty{13.6e12}{mol} adalah 0.141 hingga 0.163 miliar ton karbon setahun; dengan
\qty{44.009}{g.mol^{-1}}, 0.515 hingga 0.599 miliar ton karbon dioksida.
(b) $4.08/13.6 = \qty{30}{\percent}$ hingga $4.08/11.7 = \qty{35}{\percent}$.
(c) Setengah karbon dioksida yang dikonsumsi tersimpan ketika kalsium dan magnesium
mengendap sebagai karbonat: 0.070 hingga 0.082 miliar ton karbon setahun. Angka
nyatanya lebih kecil, karena sebagian bikarbonat pergi bersama natrium dan kalium,
yang tidak membentuk karbonat di dasar laut.
\end{solution}

\begin{solution}{exo:b1:weathering-clays:12}
Si/Al $= (w_{\ce{SiO2}}/60.08)/(2w_{\ce{Al2O3}}/101.96)$. A: $0.7748/0.7748 =
1.00$, rasio \omterm{def:b1:weathering-clays:clay-mineral}{kaolinit}, \omterm{def:b1:weathering-clays:hydrolysis-degree}{monosiallitisasi}. B: $0.9987/0.4904 = 2.04$, sebuah \omterm{def:b1:weathering-clays:clay-mineral}{mineral
lempung} 2:1, \omterm{def:b1:weathering-clays:hydrolysis-degree}{bisiallitisasi}. C: $0.0333/1.2358 = 0.027$, hampir tidak ada
silikon yang tersisa, gibsit, \omterm{def:b1:weathering-clays:hydrolysis-degree}{alitisasi}. Dari pencucian paling sedikit hingga
paling intens, yaitu menuju iklim yang paling panas dan paling basah: B, A, C.
\end{solution}

\begin{solution}{pb:b1:weathering-clays:1}
\textbf{1.} \ce{2 NaAlSi3O8 + 2 H+ + 9 H2O -> Al2Si2O5(OH)4 + 2 Na+ + 4 H4SiO4}.

\textbf{2.} \ce{CaAl2Si2O8 + 2 H+ + H2O -> Al2Si2O5(OH)4 + Ca^2+}. Tidak ada silikon
yang keluar: anortit sudah memiliki satu silikon per aluminium seperti \omterm{def:b1:weathering-clays:clay-mineral}{kaolinit}.

\textbf{3.} Satu gibsit mengambil satu aluminium, setengah anortit: dikalikan dua,
\ce{CaAl2Si2O8 + 2 H+ + 6 H2O -> 2 Al(OH)3 + Ca^2+ + 2 H4SiO4}. Oksigen:
$8 + 6 = 6 + 8$; hidrogen: $2 + 12 = 6 + 8$.

\textbf{4.} \ce{5 NaAlSi3O8 + 4 H+ + 16 H2O -> 2 Na_{0.5}Al2(Si_{3.5}Al_{0.5})O10(OH)2 + 4 Na+ + 8 H4SiO4}; % ce-unbalanced-ok: decimal subscripts the checker cannot read; balance tested in tests/bachelor
hidrogen, $4 + 32 = 36$ di ruas kiri dan $2\times2 + 8\times4 = 36$ di
ruas kanan, seimbang.

\textbf{5.} Setiap \ce{H+} disertai satu \ce{HCO3-} dari satu \ce{CO2} dan satu
\ce{H2O}:
\ce{2 NaAlSi3O8 + 2 CO2 + 11 H2O -> Al2Si2O5(OH)4 + 2 Na+ + 2 HCO3- + 4 H4SiO4}
dan \ce{CaAl2Si2O8 + 2 CO2 + 3 H2O -> Al2Si2O5(OH)4 + Ca^2+ + 2 HCO3-}.

\textbf{6.} Albit menjadi beidelit adalah \omterm{def:b1:weathering-clays:hydrolysis-degree}{bisiallitisasi}; albit dan anortit menjadi
\omterm{def:b1:weathering-clays:clay-mineral}{kaolinit} adalah \omterm{def:b1:weathering-clays:hydrolysis-degree}{monosiallitisasi}; anortit menjadi gibsit adalah \omterm{def:b1:weathering-clays:hydrolysis-degree}{alitisasi}.

\textbf{7.} $\num{e6}/262.22 = \qty{3813.6}{mol}$ albit;
$\num{e6}/278.20 = \qty{3594.5}{mol}$ anortit.

\textbf{8.} Albit: $3813.6/2 = \qty{1906.8}{mol}$ \omterm{def:b1:weathering-clays:clay-mineral}{kaolinit},
\qty{492.3}{kg}. Anortit: \qty{3594.5}{mol}, \qty{928.0}{kg}.

\textbf{9.} Albit: $2\times3813.6 = \qty{7627.2}{mol}$ dari \ce{SiO2},
\qty{458.2}{kg}, dan \qty{3813.6}{mol} natrium, \qty{87.7}{kg}. Anortit:
tidak ada silika, dan \qty{3594.5}{mol} kalsium, \qty{144.1}{kg}.

\textbf{10.} Albit: $1/2.62 = \qty{0.382}{m^3}$; \omterm{def:b1:weathering-clays:clay-mineral}{kaolinitnya},
$0.4923/2.6 = \qty{0.189}{m^3}$: padatan kehilangan \qty{50.4}{\percent} dari
volumenya. Anortit: $1/2.75 = \qty{0.364}{m^3}$; \omterm{def:b1:weathering-clays:clay-mineral}{kaolinitnya},
$0.9280/2.6 = \qty{0.357}{m^3}$: hanya kehilangan \qty{1.8}{\percent}.

\textbf{11.} Ruang yang tersisa menjadi pori, terisi air dan udara: sebuah
saprolit di atas granit yang kaya albit ringan, berpori dan rapuh, butir kuarsanya
lepas dalam sedikit lempung. Anortit berubah menjadi \omterm{def:b1:weathering-clays:clay-mineral}{kaolinit} yang hampir sama
volumenya dengannya, sehingga gabro lapuk mempertahankan sebagian besar materi padatnya.

\textbf{12.} Natrium ($z/r = 0.98$) dan kalsium (2.0) adalah kation larut dan
keluar; aluminium (5.6) adalah hidrolisat dan tinggal seluruhnya dalam \omterm{def:b1:weathering-clays:clay-mineral}{kaolinit};
silikon (10.0), pada batas, keluar sebagian dari albit dan tidak sama sekali
dari anortit, yang silikonnya dapat ditampung \omterm{def:b1:weathering-clays:clay-mineral}{kaolinit}.

\textbf{13.} Satu \ce{CO2} per albit: \qty{3813.6}{mol}, $3813.6 \times 44.01 =
\qty{167.8}{kg}$.

\textbf{14.} Dua per anortit: \qty{7189.0}{mol}, \qty{316.4}{kg}.

\textbf{15.} $7189.0 \times 61.02 = \qty{438.7}{kg}$ bikarbonat.

\textbf{16.} Organisme mengendapkannya sebagai kalsium karbonat,
\ce{Ca^2+ + 2 HCO3- -> CaCO3 + CO2 + H2O}: untuk setiap kalsium, satu dari dua
\ce{CO2} kembali ke udara, \qty{3594.5}{mol} atau \qty{158.2}{kg}, dan yang
lain terkubur dalam karbonat.

\textbf{17.} Natrium tidak membentuk karbonat di dasar laut: bikarbonatnya tidak
dapat tersimpan sebagai batuan karbonat, dan penyimpanan jangka panjang dihitung
dari kalsium dan magnesium saja.

\textbf{18.} Satu ton kalsit mengonsumsi \qty{9991}{mol} \ce{CO2} ketika melarut
dan mengembalikan semuanya ketika mengendap: tidak ada yang tersimpan, dan hanya
pelapukan silikat kalsium dan magnesium yang menarik karbon dioksida turun untuk selamanya.

\textbf{19.} $\qty{200}{m^3.s^{-1}} \times \qty{0.3}{mol.m^{-3}} =
\qty{60}{mol.s^{-1}}$; selama setahun \qty{3.156e7}{s}, \qty{1.893e9}{mol}, atau
$\num{1.893e9}/\num{50000} = \qty{37870}{mol}$ per kilometer persegi.

\textbf{20.} Dua bikarbonat per anortit: \qty{9.467e8}{mol} setahun,
$\num{9.467e8}\times\qty{278.20}{g} = \qty{263400}{t}$ setahun, atau
\qty{5.27}{t} per kilometer persegi.

\textbf{21.} Batuan segar yang terkonsumsi adalah $5.27/0.20 = \qty{26.3}{t}$, atau
$\num{26340}/2700 = \qty{9.75}{m^3}$, per kilometer persegi dan tahun: lapisan
setebal \qty{9.75}{\micro m} setahun, \qty{9.8}{m} per juta tahun.

\textbf{22.} Satu \ce{CO2} tersimpan per kalsium: $\num{9.467e8} \times 44.01 =
\qty{4.17e10}{g}$, sekitar \qty{41700}{t} setahun.

\textbf{23.} Ini hanya menghitung materi yang terbawa dalam larutan. \omterm{def:b1:weathering-clays:clay-mineral}{Kaolinit} dan
kuarsa yang ditinggalkan oleh pelapukan juga dibawa pergi oleh erosi, sebagai
lumpur dan pasir dalam sungai, dan tidak dihitung; begitu pula pelapukan \omterm{def:b1:rock-forming-minerals:mineral}{mineral}
lain dari batuan itu. Dan jika sebagian bikarbonat berasal dari \omterm{def:b1:rock-forming-minerals:groups}{feldspar} natrium
atau kalium, yang memberi satu bikarbonat per satuan rumus alih-alih dua, lebih
banyak batuan harus lapuk untuk memasoknya.

\textbf{24.} Menurut \cref{prop:b1:weathering-clays:co2}, satu mol \ce{CO2} tersimpan
per mol anortit: $3594.5 \times \qty{44.01}{g} = \qty{158.19}{kg}$.
Pelapukan satu ton anortit menarik turun untuk selamanya
\textbf{158}~kg karbon dioksida.
\end{solution}
